\hypertarget{anexo-a.-clasificaciuxf3n-de-los-ejercicios}{%
\section{Anexo A. Clasificación de los ejercicios}\label{anexo-a.-clasificaciuxf3n-de-los-ejercicios}}

\begin{longtable}[]{@{}
  >{\raggedright\arraybackslash}p{(\columnwidth - 4\tabcolsep) * \real{0.1875}}
  >{\raggedright\arraybackslash}p{(\columnwidth - 4\tabcolsep) * \real{0.1875}}
  >{\raggedright\arraybackslash}p{(\columnwidth - 4\tabcolsep) * \real{0.6250}}@{}}
\caption{Clasificación de los ejercicios de las Listas 2 y 3.}\tabularnewline
\toprule\noalign{}
\begin{minipage}[b]{\linewidth}\raggedright
Ejercicio
\end{minipage} & \begin{minipage}[b]{\linewidth}\raggedright
Sección
\end{minipage} & \begin{minipage}[b]{\linewidth}\raggedright
Justificación
\end{minipage} \\
\midrule\noalign{}
\endfirsthead
\toprule\noalign{}
\begin{minipage}[b]{\linewidth}\raggedright
Ejercicio
\end{minipage} & \begin{minipage}[b]{\linewidth}\raggedright
Sección
\end{minipage} & \begin{minipage}[b]{\linewidth}\raggedright
Justificación
\end{minipage} \\
\midrule\noalign{}
\endhead
\bottomrule\noalign{}
\endlastfoot
Lista 2 · Ejercicio 094 & AFD & Se decide con información acotada del prefijo: basta un AFD. \\
Lista 2 · Ejercicio 095 & AFD & Se decide con información acotada del prefijo: basta un AFD. \\
Lista 2 · Ejercicio 096 & AFD & Se decide con información acotada del prefijo: basta un AFD. \\
Lista 2 · Ejercicio 097 & AFD & Se decide con información acotada del prefijo: basta un AFD. \\
Lista 2 · Ejercicio 098 & AFD & Se decide con información acotada del prefijo: basta un AFD. \\
Lista 2 · Ejercicio 099 & AFD & Se decide con información acotada del prefijo: basta un AFD. \\
Lista 2 · Ejercicio 100 & AFD & Se decide con información acotada del prefijo: basta un AFD. \\
Lista 2 · Ejercicio 101 & AFD & Se decide con información acotada del prefijo: basta un AFD. \\
Lista 2 · Ejercicio 102 & AFD & Se decide con información acotada del prefijo: basta un AFD. \\
Lista 2 · Ejercicio 103 & AFD & Se decide con información acotada del prefijo: basta un AFD. \\
Lista 2 · Ejercicio 106 & AFD & Se decide con información acotada del prefijo: basta un AFD. \\
Lista 2 · Ejercicio 104 & AFD & Se decide con información acotada del prefijo: basta un AFD. \\
Lista 2 · Ejercicio 105 & AFD & Se decide con información acotada del prefijo: basta un AFD. \\
Lista 2 · Ejercicio 107 & AFD & Se decide con información acotada del prefijo: basta un AFD. \\
Lista 2 · Ejercicio 109 & AFD & Se decide con información acotada del prefijo: basta un AFD. \\
Lista 2 · Ejercicio 108 & AFD & Se decide con información acotada del prefijo: basta un AFD. \\
Lista 2 · Ejercicio 110 & AFD & Se decide con información acotada del prefijo: basta un AFD. \\
Lista 2 · Ejercicio 111 & AFD & Se decide con información acotada del prefijo: basta un AFD. \\
Lista 2 · Ejercicio 112 & AFD & Se decide con información acotada del prefijo: basta un AFD. \\
Lista 2 · Ejercicio 113 & AFD & Se decide con información acotada del prefijo: basta un AFD. \\
Lista 2 · Ejercicio 114 & AFD & Se decide con información acotada del prefijo: basta un AFD. \\
Lista 2 · Ejercicio 115 & AFD & Se decide con información acotada del prefijo: basta un AFD. \\
Lista 2 · Ejercicio 116 & --- & No regular: exige comparar el número de aes con el de ces (lema de bombeo). \\
Lista 2 · Ejercicio 117 & AFD & Se decide con información acotada del prefijo: basta un AFD. \\
Lista 2 · Ejercicio 118 & AFD & Se decide con información acotada del prefijo: basta un AFD. \\
Lista 2 · Ejercicio 119 & AFD & Se decide con información acotada del prefijo: basta un AFD. \\
Lista 2 · Ejercicio 120 & AFD & Se decide con información acotada del prefijo: basta un AFD. \\
Lista 2 · Ejercicio 121 & AFD & Se decide con información acotada del prefijo: basta un AFD. \\
Lista 2 · Ejercicio 122 & AFD & Se decide con información acotada del prefijo: basta un AFD. \\
Lista 2 · Ejercicio 123 & AFD & Se decide con información acotada del prefijo: basta un AFD. \\
Lista 2 · Ejercicio 124 & AFD & Se decide con información acotada del prefijo: basta un AFD. \\
Lista 2 · Ejercicio 125 & AFD & Propiedad de los AFD completos; se ilustra con un AFD. \\
Lista 2 · Ejercicio 126 & AFD & Propiedad de los AFD; se ilustra con un AFD. \\
Lista 2 · Ejercicio 127 & AFD & Propiedad de los AFD; se ilustra con un AFD. \\
Lista 2 · Ejercicio 128 & AFD & Se decide con información acotada del prefijo: basta un AFD. \\
Lista 2 · Ejercicio 129 & AFD & El enunciado da la quíntupla de un AFD. \\
Lista 2 · Ejercicio 130 & AFD & El enunciado da la quíntupla de un AFD. \\
Lista 2 · Ejercicio 131 & AFD & El enunciado da la quíntupla de un AFD. \\
Lista 2 · Ejercicio 132 & AFD & Se decide con información acotada del prefijo: basta un AFD. \\
Lista 2 · Ejercicio 133 & AFD & El enunciado da la quíntupla de un AFD. \\
Lista 2 · Ejercicio 134 & AFD & Se decide con información acotada del prefijo: basta un AFD. \\
Lista 2 · Ejercicio 135 & AFD & Se decide con información acotada del prefijo: basta un AFD. \\
Lista 2 · Ejercicio 136 & AFD & Se decide con información acotada del prefijo: basta un AFD. \\
Lista 2 · Ejercicio 137 & AFD & Se decide con información acotada del prefijo: basta un AFD. \\
Lista 2 · Ejercicio 138 & AFD & Se decide con información acotada del prefijo: basta un AFD. \\
Lista 2 · Ejercicio 139 & AFD & Se decide con información acotada del prefijo: basta un AFD. \\
Lista 2 · Ejercicio 140a & AFD & Se decide con información acotada del prefijo: basta un AFD. \\
Lista 2 · Ejercicio 140b & AFD & Se decide con información acotada del prefijo: basta un AFD. \\
Lista 2 · Ejercicio 141 & AFD & Se decide con información acotada del prefijo: basta un AFD. \\
Lista 2 · Ejercicio 142a & AFD & Se decide con información acotada del prefijo: basta un AFD. \\
Lista 2 · Ejercicio 142b & AFD & Se decide con información acotada del prefijo: basta un AFD. \\
Lista 2 · Ejercicio 143 & --- & Sin figuras en el enunciado. \\
Lista 2 · Ejercicio 144a & AFD & Se decide con información acotada del prefijo: basta un AFD. \\
Lista 2 · Ejercicio 144b & AFD & Se decide con información acotada del prefijo: basta un AFD. \\
Lista 2 · Ejercicio 144c & AFD & Se decide con información acotada del prefijo: basta un AFD. \\
Lista 2 · Ejercicio 144d & AFD & Se decide con información acotada del prefijo: basta un AFD. \\
Lista 2 · Ejercicio 144e & AFD & Se decide con información acotada del prefijo: basta un AFD. \\
Lista 2 · Ejercicio 145 & AFD & Se decide con información acotada del prefijo: basta un AFD. \\
Lista 2 · Ejercicio 146 & AFD & Se decide con información acotada del prefijo: basta un AFD. \\
Lista 2 · Ejercicio 147 & AFD & Se decide con información acotada del prefijo: basta un AFD. \\
Lista 2 · Ejercicio 148 & AFD & Se decide con información acotada del prefijo: basta un AFD. \\
Lista 2 · Ejercicio 149 & AFD & El enunciado da la tabla de un AFD. \\
Lista 2 · Ejercicio 150 & AFD & El enunciado da la tabla de un AFD. \\
Lista 3 · Ejercicio 151 & AFD & El enunciado da la tabla de un AFD. \\
Lista 3 · Ejercicio 152 & AFD & El enunciado da la tabla de un AFD. \\
Lista 3 · Ejercicio 153 & AFD & El enunciado da la tabla de un AFD. \\
Lista 3 · Ejercicio 154 & AFD & El enunciado da la tabla de un AFD. \\
Lista 3 · Ejercicio 155 & AFD & El enunciado da la tabla de un AFD. \\
Lista 3 · Ejercicio 156 & AFD & El enunciado da la tabla de un AFD. \\
Lista 3 · Ejercicio 157 & AFD & El enunciado da la tabla de un AFD. \\
Lista 3 · Ejercicio 158 & AFD & Función de transición determinista (parcial), aunque aparezca en la sección de AFN. \\
Lista 3 · Ejercicio 159 & AFD & Función de transición determinista (parcial), aunque aparezca en la sección de AFN. \\
Lista 3 · Ejercicio 160 & AFD & Función de transición determinista (parcial), aunque aparezca en la sección de AFN. \\
Lista 3 · Ejercicio 161 & AFD & Función de transición determinista (parcial), aunque aparezca en la sección de AFN. \\
Lista 3 · AFN del problema 094 & AFND & Adivina una posición o una rama: AFN. \\
Lista 3 · AFN del problema 095 & AFD & Pide un AFN; sin una versión no determinista más simple, se usa el AFD parcial. \\
Lista 3 · AFN del problema 096 & AFD & Pide un AFN; sin una versión no determinista más simple, se usa el AFD parcial. \\
Lista 3 · AFN del problema 097 & AFD & Pide un AFN; sin una versión no determinista más simple, se usa el AFD parcial. \\
Lista 3 · AFN del problema 098 & AFND & Adivina una posición o una rama: AFN. \\
Lista 3 · AFN del problema 099 & AFND & Adivina una posición o una rama: AFN. \\
Lista 3 · AFN del problema 100 & AFD & Pide un AFN; sin una versión no determinista más simple, se usa el AFD parcial. \\
Lista 3 · AFN del problema 101 & AFND & Adivina una posición o una rama: AFN. \\
Lista 3 · AFN del problema 102 & AFND & Adivina una posición o una rama: AFN. \\
Lista 3 · AFN del problema 103 & AFND & Adivina una posición o una rama: AFN. \\
Lista 3 · AFN del problema 104 & AFD & Pide un AFN; sin una versión no determinista más simple, se usa el AFD parcial. \\
Lista 3 · AFN del problema 105 & AFND & Adivina una posición o una rama: AFN. \\
Lista 3 · AFN del problema 106 & AFND & Adivina una posición o una rama: AFN. \\
Lista 3 · AFN del problema 107 & AFND & Adivina una posición o una rama: AFN. \\
Lista 3 · AFN del problema 108 & AFND & Adivina una posición o una rama: AFN. \\
Lista 3 · AFN del problema 109 & AFND & Adivina una posición o una rama: AFN. \\
Lista 3 · AFN del problema 110 & AFD & Pide un AFN; sin una versión no determinista más simple, se usa el AFD parcial. \\
Lista 3 · AFN del problema 111 & AFD & Pide un AFN; sin una versión no determinista más simple, se usa el AFD parcial. \\
Lista 3 · AFN del problema 112 & AFND & Adivina una posición o una rama: AFN. \\
Lista 3 · AFN del problema 113 & AFND & Adivina una posición o una rama: AFN. \\
Lista 3 · AFN del problema 114 & AFD & Pide un AFN; sin una versión no determinista más simple, se usa el AFD parcial. \\
Lista 3 · AFN del problema 115 & AFD & Pide un AFN; sin una versión no determinista más simple, se usa el AFD parcial. \\
Lista 3 · AFN del problema 116 & --- & No regular. \\
Lista 3 · AFN del problema 117 & AFD & Pide un AFN; sin una versión no determinista más simple, se usa el AFD parcial. \\
Lista 3 · AFN del problema 118 & AFND & Adivina una posición o una rama: AFN. \\
Lista 3 · AFN del problema 119 & AFD & Pide un AFN; sin una versión no determinista más simple, se usa el AFD parcial. \\
Lista 3 · AFN del problema 120 & AFND & Adivina una posición o una rama: AFN. \\
Lista 3 · AFN del problema 121 & AFD & Pide un AFN; sin una versión no determinista más simple, se usa el AFD parcial. \\
Lista 3 · AFN del problema 122 & AFD & Pide un AFN; sin una versión no determinista más simple, se usa el AFD parcial. \\
Lista 3 · AFN del problema 123 & AFND & Adivina una posición o una rama: AFN. \\
Lista 3 · Ejercicio 189 & AFND & Unión de dos búsquedas: AFN. \\
Lista 3 · Ejercicio 190 & AFND & Intersección de dos AFN. \\
Lista 3 · Ejercicio 191 & --- & Encabezado de los ejercicios 192-195. \\
Lista 3 · Ejercicio 192 & AFND & Hay pares (estado, símbolo) con varios destinos. \\
Lista 3 · Ejercicio 193 & AFND & Hay pares (estado, símbolo) con varios destinos. \\
Lista 3 · Ejercicio 194 & AFND & Hay pares (estado, símbolo) con varios destinos. \\
Lista 3 · Ejercicio 195 & AFND & Hay pares (estado, símbolo) con varios destinos. \\
Lista 3 · Ejercicio 196 & AFND-λ & La tabla tiene transiciones λ. \\
Lista 3 · Ejercicio 197 & AFND-λ & La tabla tiene transiciones λ. \\
Lista 3 · Ejercicio 198 & AFND-λ & La tabla tiene transiciones λ. \\
Lista 3 · Ejercicio 199 & AFND-λ & La tabla tiene transiciones λ. \\
Lista 3 · Ejercicio 200 & AFND-λ & La tabla tiene transiciones λ. \\
Lista 3 · Ejercicio 201-235 & --- & Pertenencia a expresiones regulares (AFN-λ de Thompson). \\
Lista 3 · Ejercicio 236 & AFND-λ & Se parte de una expresión regular: su autómata natural es el AFN-λ de Thompson. \\
Lista 3 · Ejercicio 237 & AFND-λ & Se parte de una expresión regular: su autómata natural es el AFN-λ de Thompson. \\
Lista 3 · Ejercicio 238 & AFND-λ & Se parte de una expresión regular: su autómata natural es el AFN-λ de Thompson. \\
Lista 3 · Ejercicio 239 & AFND-λ & Se parte de una expresión regular: su autómata natural es el AFN-λ de Thompson. \\
Lista 3 · Ejercicio 240 & AFND-λ & Se parte de una expresión regular: su autómata natural es el AFN-λ de Thompson. \\
Lista 3 · Ejercicio 241 & AFND-λ & Expresión regular → AFN-λ. \\
Lista 3 · Ejercicio 242 & AFND-λ & Expresión regular → AFN-λ. \\
Lista 3 · Ejercicio 243 & AFND-λ & Expresión regular → AFN-λ. \\
Lista 3 · Ejercicio 244 & AFND-λ & Expresión regular → AFN-λ. \\
Lista 3 · Ejercicio 245 & AFND-λ & Expresión regular → AFN-λ. \\
Lista 3 · Ejercicio 246 & AFND-λ & Expresión regular → AFN-λ. \\
Lista 3 · Ejercicio 247 & AFND-λ & Expresión regular → AFN-λ. \\
Lista 3 · Ejercicio 248 & AFND-λ & Expresión regular → AFN-λ. \\
Lista 3 · Ejercicio 249 & AFND-λ & Expresión regular → AFN-λ. \\
Lista 3 · Ejercicio 250 & AFND-λ & Expresión regular → AFN-λ. \\
Lista 3 · Ejercicio 251 & AFND-λ & Expresión regular → AFN-λ. \\
Lista 3 · Ejercicio 252 & AFND-λ & Expresión regular → AFN-λ. \\
Lista 3 · Ejercicio 253 & AFND-λ & Expresión regular → AFN-λ. \\
Lista 3 · Ejercicio 254 & AFD & El enunciado da la tabla de un AFD. \\
Lista 3 · Ejercicio 255 & AFD & El enunciado da la tabla de un AFD. \\
Lista 3 · Ejercicio 256 & AFD & El enunciado da la tabla de un AFD. \\
Lista 3 · Ejercicio 257 & AFD & El enunciado da la tabla de un AFD. \\
Lista 3 · Ejercicio 258 & AFD & El enunciado da la tabla de un AFD. \\
Lista 3 · Ejercicio 259 & AFD & El enunciado da la tabla de un AFD. \\
Lista 3 · Ejercicio 260 & AFND & AFN con varios destinos por par (estado, símbolo). \\
Lista 3 · Ejercicio 261 & AFND & AFN con varios destinos por par (estado, símbolo). \\
Lista 3 · Ejercicio 262 & AFD & Diseño natural con dos ramas (AFN). \\
Lista 3 · Ejercicio 263 & AFND & AFN con varios destinos por par (estado, símbolo). \\
\end{longtable}

\hypertarget{anexo-b.-cerraduras-ux3bb}{%
\section{Anexo B. Cerraduras λ}\label{anexo-b.-cerraduras-ux3bb}}

Para cada AFN-λ se indica por qué hacen falta las transiciones λ y la cerradura de cada estado. La cerradura es lo primero que calcula el simulador: el inicio de la lectura no es \{q₀\} sino λ-cl(q₀), y después de cada símbolo se vuelve a cerrar el conjunto de estados activos.

\textbf{Lista 3 · Ejercicio 196.} \textbf{Cerraduras:} Cl(A) = \{A, B\}, Cl(B) = \{B\}, Cl(C) = \{C\}. \textbf{a)} δ̂(A, 110) = ∅. \textbf{b)} Cl(\{A\}) = \{A, B, C\}. \textbf{c)} Cl(\{B\}) = \{B, C\}. \textbf{d)} Cl(\{A, B\}) = \{A, B, C\}. \textbf{e)} δ(A, 0) = \{A\} (sin cerradura). \textbf{AFD equivalente:} 3 estados.

λ-cl(A) = \{A, B\} · λ-cl(B) = \{B\} · λ-cl(C) = \{C\}

\textbf{Lista 3 · Ejercicio 197.} \textbf{Cerraduras:} Cl(A) = \{A, B, C, D\}, Cl(B) = \{B, D\}, Cl(C) = \{B, C, D\}, Cl(D) = \{D\}. \textbf{a)} Cl(\{A\}) = \{A, B, C, D\}. \textbf{b)} Cl(\{B, D\}) = \{B, C, D\}. \textbf{c)} Cl(\{A, B, C\}) = \{A, B, C, D\}. \textbf{d)} δ(A, 0) = \{B, D\} (sin cerradura). \textbf{e)} δ̂(\{A\}, 010) = \{A, B, C, D\}. \textbf{AFD equivalente:} 1 estados.

λ-cl(A) = \{A, B, C, D\} · λ-cl(B) = \{B, D\} · λ-cl(C) = \{B, C, D\} · λ-cl(D) = \{D\}

\textbf{Lista 3 · Ejercicio 198.} \textbf{Cerraduras:} Cl(A) = \{A, B, C\}, Cl(B) = \{B\}, Cl(C) = \{B, C\}, Cl(D) = \{D\}. \textbf{a)} Cl(\{A\}) = \{A, B, C\}. \textbf{b)} Cl(\{C, D\}) = \{B, C, D\}. \textbf{c)} Cl(\{A, C, D\}) = \{A, B, C, D\}. \textbf{d)} δ̂(B, 0111) = \{B, C, D\}. \textbf{AFD equivalente:} 3 estados.

λ-cl(A) = \{A, B, C\} · λ-cl(B) = \{B\} · λ-cl(C) = \{B, C\} · λ-cl(D) = \{D\}

\textbf{Lista 3 · Ejercicio 199.} \textbf{Cerraduras:} Cl(A) = \{A, B, C, E\}, Cl(B) = \{B\}, Cl(C) = \{C, E\}, Cl(D) = \{D\}, Cl(E) = \{E\}. \textbf{a)} Cl(\{A\}) = \{A, B, C, E\}. \textbf{b)} Cl(\{C\}) = \{C, E\}. \textbf{c)} δ̂(A, 011) = \{D\}. \textbf{AFD equivalente:} 5 estados.

λ-cl(A) = \{A, B, C, E\} · λ-cl(B) = \{B\} · λ-cl(C) = \{C, E\} · λ-cl(D) = \{D\} · λ-cl(E) = \{E\}

\textbf{Lista 3 · Ejercicio 200.} \textbf{Cerraduras:} Cl(A) = \{A, B\}, Cl(B) = \{B\}, Cl(C) = \{C, E\}, Cl(D) = \{D, E, F\}, Cl(E) = \{E\}, Cl(F) = \{E, F\}. \textbf{a)} Cl(\{A\}) = \{A, B, C, E\}. \textbf{b)} Cl(\{C, D\}) = \{C, D, E, F\}. \textbf{c)} δ̂(C, 1011) = \{A, B, C, D, E, F\}. \textbf{AFD equivalente:} 6 estados.

λ-cl(A) = \{A, B\} · λ-cl(B) = \{B\} · λ-cl(C) = \{C, E\} · λ-cl(D) = \{D, E, F\} · λ-cl(E) = \{E\} · λ-cl(F) = \{E, F\}

\textbf{Lista 3 · Ejercicio 236.} \textbf{Expresión regular:} \texttt{(0\ +\ 1)*001(0\ +\ 1)*}.

λ-cl(q6) = \{q6, q4, q0, q2, q7, q8\} · λ-cl(q4) = \{q4, q0, q2\} · λ-cl(q0) = \{q0\} · λ-cl(q1) = \{q4, q0, q1, q2, q5, q7, q8\} · λ-cl(q2) = \{q2\} · λ-cl(q3) = \{q4, q0, q2, q3, q5, q7, q8\} · λ-cl(q5) = \{q4, q0, q2, q5, q7, q8\} · λ-cl(q7) = \{q7, q8\} · λ-cl(q8) = \{q8\} · λ-cl(q9) = \{q9, q10\} · λ-cl(q10) = \{q10\} · λ-cl(q11) = \{q11, q12\} · λ-cl(q12) = \{q12\} · λ-cl(q13) = \{q13, q20, q18, q14, q16, q21\} · λ-cl(q20) = \{q20, q18, q14, q16, q21\} · λ-cl(q18) = \{q18, q14, q16\} · λ-cl(q14) = \{q14\} · λ-cl(q15) = \{q18, q14, q15, q16, q19, q21\} · λ-cl(q16) = \{q16\} · λ-cl(q17) = \{q18, q14, q16, q17, q19, q21\} · λ-cl(q19) = \{q18, q14, q16, q19, q21\} · λ-cl(q21) = \{q21\}

\textbf{Lista 3 · Ejercicio 237.} \textbf{Expresión regular:} \texttt{(0\ +\ 1)*010}.

λ-cl(q6) = \{q6, q4, q0, q2, q7, q8\} · λ-cl(q4) = \{q4, q0, q2\} · λ-cl(q0) = \{q0\} · λ-cl(q1) = \{q4, q0, q1, q2, q5, q7, q8\} · λ-cl(q2) = \{q2\} · λ-cl(q3) = \{q4, q0, q2, q3, q5, q7, q8\} · λ-cl(q5) = \{q4, q0, q2, q5, q7, q8\} · λ-cl(q7) = \{q7, q8\} · λ-cl(q8) = \{q8\} · λ-cl(q9) = \{q9, q10\} · λ-cl(q10) = \{q10\} · λ-cl(q11) = \{q11, q12\} · λ-cl(q12) = \{q12\} · λ-cl(q13) = \{q13\}

\textbf{Lista 3 · Ejercicio 238.} \textbf{Expresión regular:} \texttt{(0\ +\ 1)(0\ +\ 1)(0\ +\ 1)(λ\ +\ (0\ +\ 1)(0\ +\ 1))}.

λ-cl(q4) = \{q4, q0, q2\} · λ-cl(q0) = \{q0\} · λ-cl(q1) = \{q1, q5, q10, q6, q8\} · λ-cl(q2) = \{q2\} · λ-cl(q3) = \{q3, q5, q10, q6, q8\} · λ-cl(q5) = \{q5, q10, q6, q8\} · λ-cl(q10) = \{q10, q6, q8\} · λ-cl(q6) = \{q6\} · λ-cl(q7) = \{q7, q11, q16, q12, q14\} · λ-cl(q8) = \{q8\} · λ-cl(q9) = \{q9, q11, q16, q12, q14\} · λ-cl(q11) = \{q11, q16, q12, q14\} · λ-cl(q16) = \{q16, q12, q14\} · λ-cl(q12) = \{q12\} · λ-cl(q13) = \{q13, q17, q32, q18, q19, q24, q20, q22, q33\} · λ-cl(q14) = \{q14\} · λ-cl(q15) = \{q15, q17, q32, q18, q19, q24, q20, q22, q33\} · λ-cl(q17) = \{q17, q32, q18, q19, q24, q20, q22, q33\} · λ-cl(q32) = \{q32, q18, q19, q24, q20, q22, q33\} · λ-cl(q18) = \{q18, q19, q33\} · λ-cl(q19) = \{q19, q33\} · λ-cl(q24) = \{q24, q20, q22\} · λ-cl(q20) = \{q20\} · λ-cl(q21) = \{q21, q25, q30, q26, q28\} · λ-cl(q22) = \{q22\} · λ-cl(q23) = \{q23, q25, q30, q26, q28\} · λ-cl(q25) = \{q25, q30, q26, q28\} · λ-cl(q30) = \{q30, q26, q28\} · λ-cl(q26) = \{q26\} · λ-cl(q27) = \{q27, q31, q33\} · λ-cl(q28) = \{q28\} · λ-cl(q29) = \{q29, q31, q33\} · λ-cl(q31) = \{q31, q33\} · λ-cl(q33) = \{q33\}

\textbf{Lista 3 · Ejercicio 239.} \textbf{Expresión regular:} \texttt{01(0\ +\ 1)*\ +\ (0\ +\ 1)*10}.

λ-cl(q24) = \{q24, q0, q18, q16, q12, q14, q19, q20\} · λ-cl(q0) = \{q0\} · λ-cl(q1) = \{q1, q2\} · λ-cl(q2) = \{q2\} · λ-cl(q3) = \{q3, q10, q8, q4, q6, q11, q25\} · λ-cl(q10) = \{q10, q8, q4, q6, q11, q25\} · λ-cl(q8) = \{q8, q4, q6\} · λ-cl(q4) = \{q4\} · λ-cl(q5) = \{q8, q4, q5, q6, q9, q11, q25\} · λ-cl(q6) = \{q6\} · λ-cl(q7) = \{q8, q4, q6, q7, q9, q11, q25\} · λ-cl(q9) = \{q8, q4, q6, q9, q11, q25\} · λ-cl(q11) = \{q11, q25\} · λ-cl(q18) = \{q18, q16, q12, q14, q19, q20\} · λ-cl(q16) = \{q16, q12, q14\} · λ-cl(q12) = \{q12\} · λ-cl(q13) = \{q16, q12, q13, q14, q17, q19, q20\} · λ-cl(q14) = \{q14\} · λ-cl(q15) = \{q16, q12, q14, q15, q17, q19, q20\} · λ-cl(q17) = \{q16, q12, q14, q17, q19, q20\} · λ-cl(q19) = \{q19, q20\} · λ-cl(q20) = \{q20\} · λ-cl(q21) = \{q21, q22\} · λ-cl(q22) = \{q22\} · λ-cl(q23) = \{q23, q25\} · λ-cl(q25) = \{q25\}

\textbf{Lista 3 · Ejercicio 240.} \textbf{Expresión regular:} \texttt{((0\ +\ 1)(0\ +\ 1))*01((0\ +\ 1)(0\ +\ 1))*\ +\ (0\ +\ 1)((0\ +\ 1)(0\ +\ 1))*01((0\ +\ 1)(0\ +\ 1))*(0\ +\ 1)}. Si w = x01y tiene longitud par, \textbar x\textbar{} + \textbar y\textbar{} es par, así que \textbar x\textbar{} y \textbar y\textbar{} son ambos pares o ambos impares; cada término cubre uno de los dos casos.

λ-cl(q76) = \{q76, q12, q4, q0, q2, q13, q14, q36, q32, q34\} · λ-cl(q12) = \{q12, q4, q0, q2, q13, q14\} · λ-cl(q4) = \{q4, q0, q2\} · λ-cl(q0) = \{q0\} · λ-cl(q1) = \{q1, q5, q10, q6, q8\} · λ-cl(q2) = \{q2\} · λ-cl(q3) = \{q3, q5, q10, q6, q8\} · λ-cl(q5) = \{q5, q10, q6, q8\} · λ-cl(q10) = \{q10, q6, q8\} · λ-cl(q6) = \{q6\} · λ-cl(q7) = \{q4, q0, q2, q7, q11, q13, q14\} · λ-cl(q8) = \{q8\} · λ-cl(q9) = \{q4, q0, q2, q9, q11, q13, q14\} · λ-cl(q11) = \{q4, q0, q2, q11, q13, q14\} · λ-cl(q13) = \{q13, q14\} · λ-cl(q14) = \{q14\} · λ-cl(q15) = \{q15, q16\} · λ-cl(q16) = \{q16\} · λ-cl(q17) = \{q17, q30, q22, q18, q20, q31, q77\} · λ-cl(q30) = \{q30, q22, q18, q20, q31, q77\} · λ-cl(q22) = \{q22, q18, q20\} · λ-cl(q18) = \{q18\} · λ-cl(q19) = \{q19, q23, q28, q24, q26\} · λ-cl(q20) = \{q20\} · λ-cl(q21) = \{q21, q23, q28, q24, q26\} · λ-cl(q23) = \{q23, q28, q24, q26\} · λ-cl(q28) = \{q28, q24, q26\} · λ-cl(q24) = \{q24\} · λ-cl(q25) = \{q22, q18, q20, q25, q29, q31, q77\} · λ-cl(q26) = \{q26\} · λ-cl(q27) = \{q22, q18, q20, q27, q29, q31, q77\} · λ-cl(q29) = \{q22, q18, q20, q29, q31, q77\} · λ-cl(q31) = \{q31, q77\} · λ-cl(q36) = \{q36, q32, q34\} · λ-cl(q32) = \{q32\} · λ-cl(q33) = \{q33, q37, q50, q42, q38, q40, q51, q52\} · λ-cl(q34) = \{q34\} · λ-cl(q35) = \{q35, q37, q50, q42, q38, q40, q51, q52\} · λ-cl(q37) = \{q37, q50, q42, q38, q40, q51, q52\} · λ-cl(q50) = \{q50, q42, q38, q40, q51, q52\} · λ-cl(q42) = \{q42, q38, q40\} · λ-cl(q38) = \{q38\} · λ-cl(q39) = \{q39, q43, q48, q44, q46\} · λ-cl(q40) = \{q40\} · λ-cl(q41) = \{q41, q43, q48, q44, q46\} · λ-cl(q43) = \{q43, q48, q44, q46\} · λ-cl(q48) = \{q48, q44, q46\} · λ-cl(q44) = \{q44\} · λ-cl(q45) = \{q42, q38, q40, q45, q49, q51, q52\} · λ-cl(q46) = \{q46\} · λ-cl(q47) = \{q42, q38, q40, q47, q49, q51, q52\} · λ-cl(q49) = \{q42, q38, q40, q49, q51, q52\} · λ-cl(q51) = \{q51, q52\} · λ-cl(q52) = \{q52\} · λ-cl(q53) = \{q53, q54\} · λ-cl(q54) = \{q54\} · λ-cl(q55) = \{q55, q68, q60, q56, q58, q69, q74, q70, q72\} · λ-cl(q68) = \{q68, q60, q56, q58, q69, q74, q70, q72\} · λ-cl(q60) = \{q60, q56, q58\} · λ-cl(q56) = \{q56\} · λ-cl(q57) = \{q57, q61, q66, q62, q64\} · λ-cl(q58) = \{q58\} · λ-cl(q59) = \{q59, q61, q66, q62, q64\} · λ-cl(q61) = \{q61, q66, q62, q64\} · λ-cl(q66) = \{q66, q62, q64\} · λ-cl(q62) = \{q62\} · λ-cl(q63) = \{q60, q56, q58, q63, q67, q69, q74, q70, q72\} · λ-cl(q64) = \{q64\} · λ-cl(q65) = \{q60, q56, q58, q65, q67, q69, q74, q70, q72\} · λ-cl(q67) = \{q60, q56, q58, q67, q69, q74, q70, q72\} · λ-cl(q69) = \{q69, q74, q70, q72\} · λ-cl(q74) = \{q74, q70, q72\} · λ-cl(q70) = \{q70\} · λ-cl(q71) = \{q71, q75, q77\} · λ-cl(q72) = \{q72\} · λ-cl(q73) = \{q73, q75, q77\} · λ-cl(q75) = \{q75, q77\} · λ-cl(q77) = \{q77\}

\textbf{Lista 3 · Ejercicio 241.} AFN-λ de 6 estados; AFD mínimo de 3 estados.

λ-cl(q4) = \{q4, q0, q2\} · λ-cl(q0) = \{q0\} · λ-cl(q1) = \{q1, q5\} · λ-cl(q2) = \{q2\} · λ-cl(q3) = \{q3, q5\} · λ-cl(q5) = \{q5\}

\textbf{Lista 3 · Ejercicio 242.} AFN-λ de 8 estados; AFD mínimo de 4 estados.

λ-cl(q6) = \{q6, q2, q0, q3, q4, q7\} · λ-cl(q2) = \{q2, q0, q3, q7\} · λ-cl(q0) = \{q0\} · λ-cl(q1) = \{q0, q1, q3, q7\} · λ-cl(q3) = \{q3, q7\} · λ-cl(q4) = \{q4\} · λ-cl(q5) = \{q5, q7\} · λ-cl(q7) = \{q7\}

\textbf{Lista 3 · Ejercicio 243.} AFN-λ de 8 estados; AFD mínimo de 1 estados.

λ-cl(q6) = \{q6, q4, q0, q2, q7\} · λ-cl(q4) = \{q4, q0, q2\} · λ-cl(q0) = \{q0\} · λ-cl(q1) = \{q4, q0, q1, q2, q5, q7\} · λ-cl(q2) = \{q2\} · λ-cl(q3) = \{q4, q0, q2, q3, q5, q7\} · λ-cl(q5) = \{q4, q0, q2, q5, q7\} · λ-cl(q7) = \{q7\}

\textbf{Lista 3 · Ejercicio 244.} AFN-λ de 8 estados; AFD mínimo de 4 estados.

λ-cl(q0) = \{q0\} · λ-cl(q1) = \{q1, q6, q2, q4\} · λ-cl(q6) = \{q6, q2, q4\} · λ-cl(q2) = \{q2\} · λ-cl(q3) = \{q3, q7\} · λ-cl(q4) = \{q4\} · λ-cl(q5) = \{q5, q7\} · λ-cl(q7) = \{q7\}

\textbf{Lista 3 · Ejercicio 245.} AFN-λ de 10 estados; AFD mínimo de 3 estados.

λ-cl(q0) = \{q0\} · λ-cl(q1) = \{q1, q8, q6, q2, q4, q9\} · λ-cl(q8) = \{q8, q6, q2, q4, q9\} · λ-cl(q6) = \{q6, q2, q4\} · λ-cl(q2) = \{q2\} · λ-cl(q3) = \{q6, q2, q3, q4, q7, q9\} · λ-cl(q4) = \{q4\} · λ-cl(q5) = \{q6, q2, q4, q5, q7, q9\} · λ-cl(q7) = \{q6, q2, q4, q7, q9\} · λ-cl(q9) = \{q9\}

\textbf{Lista 3 · Ejercicio 246.} AFN-λ de 10 estados; AFD mínimo de 4 estados.

λ-cl(q2) = \{q2, q0, q3, q8, q4, q6\} · λ-cl(q0) = \{q0\} · λ-cl(q1) = \{q0, q1, q3, q8, q4, q6\} · λ-cl(q3) = \{q3, q8, q4, q6\} · λ-cl(q8) = \{q8, q4, q6\} · λ-cl(q4) = \{q4\} · λ-cl(q5) = \{q5, q9\} · λ-cl(q6) = \{q6\} · λ-cl(q7) = \{q7, q9\} · λ-cl(q9) = \{q9\}

\textbf{Lista 3 · Ejercicio 247.} AFN-λ de 12 estados; AFD mínimo de 5 estados.

λ-cl(q10) = \{q10, q4, q0, q5, q8, q6, q9, q11\} · λ-cl(q4) = \{q4, q0, q5, q11\} · λ-cl(q0) = \{q0\} · λ-cl(q1) = \{q1, q2\} · λ-cl(q2) = \{q2\} · λ-cl(q3) = \{q0, q3, q5, q11\} · λ-cl(q5) = \{q5, q11\} · λ-cl(q8) = \{q8, q6, q9, q11\} · λ-cl(q6) = \{q6\} · λ-cl(q7) = \{q6, q7, q9, q11\} · λ-cl(q9) = \{q9, q11\} · λ-cl(q11) = \{q11\}

\textbf{Lista 3 · Ejercicio 248.} AFN-λ de 10 estados; AFD mínimo de 5 estados.

λ-cl(q6) = \{q6, q2, q0, q3, q4, q7, q8\} · λ-cl(q2) = \{q2, q0, q3, q7, q8\} · λ-cl(q0) = \{q0\} · λ-cl(q1) = \{q0, q1, q3, q7, q8\} · λ-cl(q3) = \{q3, q7, q8\} · λ-cl(q4) = \{q4\} · λ-cl(q5) = \{q5, q7, q8\} · λ-cl(q7) = \{q7, q8\} · λ-cl(q8) = \{q8\} · λ-cl(q9) = \{q9\}

\textbf{Lista 3 · Ejercicio 249.} AFN-λ de 12 estados; AFD mínimo de 3 estados.

λ-cl(q6) = \{q6, q2, q0, q3, q4, q7, q10, q8, q11\} · λ-cl(q2) = \{q2, q0, q3, q7, q10, q8, q11\} · λ-cl(q0) = \{q0\} · λ-cl(q1) = \{q0, q1, q3, q7, q10, q8, q11\} · λ-cl(q3) = \{q3, q7, q10, q8, q11\} · λ-cl(q4) = \{q4\} · λ-cl(q5) = \{q5, q7, q10, q8, q11\} · λ-cl(q7) = \{q7, q10, q8, q11\} · λ-cl(q10) = \{q10, q8, q11\} · λ-cl(q8) = \{q8\} · λ-cl(q9) = \{q8, q9, q11\} · λ-cl(q11) = \{q11\}

\textbf{Lista 3 · Ejercicio 250.} AFN-λ de 16 estados; AFD mínimo de 6 estados.

λ-cl(q14) = \{q14, q2, q0, q3, q6, q4, q7, q10, q8, q11, q12, q15\} · λ-cl(q2) = \{q2, q0, q3, q6, q4, q7, q15\} · λ-cl(q0) = \{q0\} · λ-cl(q1) = \{q0, q1, q3, q6, q4, q7, q15\} · λ-cl(q3) = \{q3, q6, q4, q7, q15\} · λ-cl(q6) = \{q6, q4, q7, q15\} · λ-cl(q4) = \{q4\} · λ-cl(q5) = \{q4, q5, q7, q15\} · λ-cl(q7) = \{q7, q15\} · λ-cl(q10) = \{q10, q8, q11, q12\} · λ-cl(q8) = \{q8\} · λ-cl(q9) = \{q8, q9, q11, q12\} · λ-cl(q11) = \{q11, q12\} · λ-cl(q12) = \{q12\} · λ-cl(q13) = \{q13, q15\} · λ-cl(q15) = \{q15\}

\textbf{Lista 3 · Ejercicio 251.} AFN-λ de 18 estados; AFD mínimo de 5 estados.

λ-cl(q16) = \{q16, q6, q0, q4, q2, q5, q7, q8, q12, q10, q13, q14\} · λ-cl(q6) = \{q6, q0, q4, q2, q5, q7, q8\} · λ-cl(q0) = \{q0\} · λ-cl(q1) = \{q1, q7, q8\} · λ-cl(q4) = \{q4, q2, q5, q7, q8\} · λ-cl(q2) = \{q2\} · λ-cl(q3) = \{q2, q3, q5, q7, q8\} · λ-cl(q5) = \{q5, q7, q8\} · λ-cl(q7) = \{q7, q8\} · λ-cl(q8) = \{q8\} · λ-cl(q9) = \{q9, q17\} · λ-cl(q12) = \{q12, q10, q13, q14\} · λ-cl(q10) = \{q10\} · λ-cl(q11) = \{q10, q11, q13, q14\} · λ-cl(q13) = \{q13, q14\} · λ-cl(q14) = \{q14\} · λ-cl(q15) = \{q15, q17\} · λ-cl(q17) = \{q17\}

\textbf{Lista 3 · Ejercicio 252.} AFN-λ de 38 estados; AFD mínimo de 8 estados.

λ-cl(q36) = \{q36, q6, q4, q0, q2, q7, q16, q8, q12, q26, q18, q22\} · λ-cl(q6) = \{q6, q4, q0, q2, q7, q16, q8, q12\} · λ-cl(q4) = \{q4, q0, q2\} · λ-cl(q0) = \{q0\} · λ-cl(q1) = \{q4, q0, q1, q2, q5, q7, q16, q8, q12\} · λ-cl(q2) = \{q2\} · λ-cl(q3) = \{q4, q0, q2, q3, q5, q7, q16, q8, q12\} · λ-cl(q5) = \{q4, q0, q2, q5, q7, q16, q8, q12\} · λ-cl(q7) = \{q7, q16, q8, q12\} · λ-cl(q16) = \{q16, q8, q12\} · λ-cl(q8) = \{q8\} · λ-cl(q9) = \{q9, q10\} · λ-cl(q10) = \{q10\} · λ-cl(q11) = \{q11, q17, q37\} · λ-cl(q12) = \{q12\} · λ-cl(q13) = \{q13, q14\} · λ-cl(q14) = \{q14\} · λ-cl(q15) = \{q15, q17, q37\} · λ-cl(q17) = \{q17, q37\} · λ-cl(q26) = \{q26, q18, q22\} · λ-cl(q18) = \{q18\} · λ-cl(q19) = \{q19, q20\} · λ-cl(q20) = \{q20\} · λ-cl(q21) = \{q21, q27, q34, q32, q28, q30, q35, q37\} · λ-cl(q22) = \{q22\} · λ-cl(q23) = \{q23, q24\} · λ-cl(q24) = \{q24\} · λ-cl(q25) = \{q25, q27, q34, q32, q28, q30, q35, q37\} · λ-cl(q27) = \{q27, q34, q32, q28, q30, q35, q37\} · λ-cl(q34) = \{q34, q32, q28, q30, q35, q37\} · λ-cl(q32) = \{q32, q28, q30\} · λ-cl(q28) = \{q28\} · λ-cl(q29) = \{q32, q28, q29, q30, q33, q35, q37\} · λ-cl(q30) = \{q30\} · λ-cl(q31) = \{q32, q28, q30, q31, q33, q35, q37\} · λ-cl(q33) = \{q32, q28, q30, q33, q35, q37\} · λ-cl(q35) = \{q35, q37\} · λ-cl(q37) = \{q37\}

\textbf{Lista 3 · Ejercicio 253.} AFN-λ de 44 estados; AFD mínimo de 18 estados.

λ-cl(q14) = \{q14, q0, q12, q10, q2, q6, q13, q15, q32, q30, q24, q16, q20, q26, q33, q42, q34, q35, q36, q43\} · λ-cl(q0) = \{q0\} · λ-cl(q1) = \{q1, q15, q32, q30, q24, q16, q20, q26, q33, q42, q34, q35, q36, q43\} · λ-cl(q12) = \{q12, q10, q2, q6, q13, q15, q32, q30, q24, q16, q20, q26, q33, q42, q34, q35, q36, q43\} · λ-cl(q10) = \{q10, q2, q6\} · λ-cl(q2) = \{q2\} · λ-cl(q3) = \{q3, q4\} · λ-cl(q4) = \{q4\} · λ-cl(q5) = \{q10, q2, q5, q6, q11, q13, q15, q32, q30, q24, q16, q20, q26, q33, q42, q34, q35, q36, q43\} · λ-cl(q6) = \{q6\} · λ-cl(q7) = \{q7, q8\} · λ-cl(q8) = \{q8\} · λ-cl(q9) = \{q10, q2, q6, q9, q11, q13, q15, q32, q30, q24, q16, q20, q26, q33, q42, q34, q35, q36, q43\} · λ-cl(q11) = \{q10, q2, q6, q11, q13, q15, q32, q30, q24, q16, q20, q26, q33, q42, q34, q35, q36, q43\} · λ-cl(q13) = \{q13, q15, q32, q30, q24, q16, q20, q26, q33, q42, q34, q35, q36, q43\} · λ-cl(q15) = \{q15, q32, q30, q24, q16, q20, q26, q33, q42, q34, q35, q36, q43\} · λ-cl(q32) = \{q32, q30, q24, q16, q20, q26, q33, q42, q34, q35, q36, q43\} · λ-cl(q30) = \{q30, q24, q16, q20, q26\} · λ-cl(q24) = \{q24, q16, q20\} · λ-cl(q16) = \{q16\} · λ-cl(q17) = \{q17, q18\} · λ-cl(q18) = \{q18\} · λ-cl(q19) = \{q30, q24, q16, q19, q20, q25, q26, q31, q33, q42, q34, q35, q36, q43\} · λ-cl(q20) = \{q20\} · λ-cl(q21) = \{q21, q22\} · λ-cl(q22) = \{q22\} · λ-cl(q23) = \{q30, q24, q16, q20, q23, q25, q26, q31, q33, q42, q34, q35, q36, q43\} · λ-cl(q25) = \{q30, q24, q16, q20, q25, q26, q31, q33, q42, q34, q35, q36, q43\} · λ-cl(q26) = \{q26\} · λ-cl(q27) = \{q27, q28\} · λ-cl(q28) = \{q28\} · λ-cl(q29) = \{q30, q24, q16, q20, q26, q29, q31, q33, q42, q34, q35, q36, q43\} · λ-cl(q31) = \{q30, q24, q16, q20, q26, q31, q33, q42, q34, q35, q36, q43\} · λ-cl(q33) = \{q33, q42, q34, q35, q36, q43\} · λ-cl(q42) = \{q42, q34, q35, q36, q43\} · λ-cl(q34) = \{q34, q35, q43\} · λ-cl(q35) = \{q35, q43\} · λ-cl(q36) = \{q36\} · λ-cl(q37) = \{q37, q38\} · λ-cl(q38) = \{q38\} · λ-cl(q39) = \{q39, q40\} · λ-cl(q40) = \{q40\} · λ-cl(q41) = \{q41, q43\} · λ-cl(q43) = \{q43\}

\hypertarget{anexo-c.-conversiones}{%
\section{Anexo C. Conversiones}\label{anexo-c.-conversiones}}

\begin{longtable}[]{@{}
  >{\raggedright\arraybackslash}p{(\columnwidth - 10\tabcolsep) * \real{0.1667}}
  >{\raggedright\arraybackslash}p{(\columnwidth - 10\tabcolsep) * \real{0.1667}}
  >{\centering\arraybackslash}p{(\columnwidth - 10\tabcolsep) * \real{0.1667}}
  >{\centering\arraybackslash}p{(\columnwidth - 10\tabcolsep) * \real{0.1667}}
  >{\centering\arraybackslash}p{(\columnwidth - 10\tabcolsep) * \real{0.1667}}
  >{\centering\arraybackslash}p{(\columnwidth - 10\tabcolsep) * \real{0.1667}}@{}}
\caption{Tamaño de los autómatas antes y después de convertir.}\tabularnewline
\toprule\noalign{}
\begin{minipage}[b]{\linewidth}\raggedright
Ejercicio
\end{minipage} & \begin{minipage}[b]{\linewidth}\raggedright
Tipo
\end{minipage} & \begin{minipage}[b]{\linewidth}\centering
Estados del AFN
\end{minipage} & \begin{minipage}[b]{\linewidth}\centering
Estados del AFD
\end{minipage} & \begin{minipage}[b]{\linewidth}\centering
Estados del AFD mínimo
\end{minipage} & \begin{minipage}[b]{\linewidth}\centering
Cota 2ⁿ
\end{minipage} \\
\midrule\noalign{}
\endfirsthead
\toprule\noalign{}
\begin{minipage}[b]{\linewidth}\raggedright
Ejercicio
\end{minipage} & \begin{minipage}[b]{\linewidth}\raggedright
Tipo
\end{minipage} & \begin{minipage}[b]{\linewidth}\centering
Estados del AFN
\end{minipage} & \begin{minipage}[b]{\linewidth}\centering
Estados del AFD
\end{minipage} & \begin{minipage}[b]{\linewidth}\centering
Estados del AFD mínimo
\end{minipage} & \begin{minipage}[b]{\linewidth}\centering
Cota 2ⁿ
\end{minipage} \\
\midrule\noalign{}
\endhead
\bottomrule\noalign{}
\endlastfoot
Lista 3 · AFN del problema 094 & AFN & 3 & 3 & 3 & 8 \\
Lista 3 · AFN del problema 098 & AFN & 6 & 6 & 6 & 64 \\
Lista 3 · AFN del problema 099 & AFN & 7 & 7 & 7 & 128 \\
Lista 3 · AFN del problema 101 & AFN & 4 & 6 & 4 & 16 \\
Lista 3 · AFN del problema 102 & AFN & 7 & 12 & 7 & 128 \\
Lista 3 · AFN del problema 103 & AFN & 7 & 12 & 7 & 128 \\
Lista 3 · AFN del problema 105 & AFN & 5 & 5 & 5 & 32 \\
Lista 3 · AFN del problema 106 & AFN & 8 & 14 & 8 & 256 \\
Lista 3 · AFN del problema 107 & AFN & 16 & 20 & 13 & 65536 \\
Lista 3 · AFN del problema 108 & AFN & 23 & 88 & 37 & 8388608 \\
Lista 3 · AFN del problema 109 & AFN & 26 & 40 & 22 & 67108864 \\
Lista 3 · AFN del problema 112 & AFN & 7 & 11 & 7 & 128 \\
Lista 3 · AFN del problema 113 & AFN & 11 & 13 & 7 & 2048 \\
Lista 3 · AFN del problema 118 & AFN & 5 & 5 & 5 & 32 \\
Lista 3 · AFN del problema 120 & AFN & 5 & 7 & 7 & 32 \\
Lista 3 · AFN del problema 123 & AFN & 14 & 15 & 15 & 16384 \\
Lista 3 · Ejercicio 189 & AFN & 7 & 12 & 7 & 128 \\
Lista 3 · Ejercicio 190 & AFN & 26 & 32 & 19 & 67108864 \\
Lista 3 · Ejercicio 192 & AFN & 3 & 5 & 5 & 8 \\
Lista 3 · Ejercicio 193 & AFN & 4 & 7 & 6 & 16 \\
Lista 3 · Ejercicio 194 & AFN & 4 & 4 & 2 & 16 \\
Lista 3 · Ejercicio 195 & AFN & 5 & 12 & 7 & 32 \\
Lista 3 · Ejercicio 196 & AFN-λ & 3 & 3 & 2 & 8 \\
Lista 3 · Ejercicio 197 & AFN-λ & 4 & 2 & 1 & 16 \\
Lista 3 · Ejercicio 198 & AFN-λ & 4 & 3 & 2 & 16 \\
Lista 3 · Ejercicio 199 & AFN-λ & 5 & 4 & 2 & 32 \\
Lista 3 · Ejercicio 200 & AFN-λ & 6 & 6 & 2 & 64 \\
Lista 3 · Ejercicio 236 & AFN-λ & 22 & 9 & 4 & 4194304 \\
Lista 3 · Ejercicio 237 & AFN-λ & 14 & 5 & 4 & 16384 \\
Lista 3 · Ejercicio 238 & AFN-λ & 34 & 11 & 6 & 17179869184 \\
Lista 3 · Ejercicio 239 & AFN-λ & 26 & 9 & 6 & 67108864 \\
Lista 3 · Ejercicio 240 & AFN-λ & 78 & 27 & 6 & 302231454903657293676544 \\
Lista 3 · Ejercicio 241 & AFN-λ & 6 & 3 & 2 & 64 \\
Lista 3 · Ejercicio 242 & AFN-λ & 8 & 3 & 3 & 256 \\
Lista 3 · Ejercicio 243 & AFN-λ & 8 & 3 & 1 & 256 \\
Lista 3 · Ejercicio 244 & AFN-λ & 8 & 4 & 3 & 256 \\
Lista 3 · Ejercicio 245 & AFN-λ & 10 & 4 & 2 & 1024 \\
Lista 3 · Ejercicio 246 & AFN-λ & 10 & 3 & 3 & 1024 \\
Lista 3 · Ejercicio 247 & AFN-λ & 12 & 4 & 4 & 4096 \\
Lista 3 · Ejercicio 248 & AFN-λ & 10 & 4 & 4 & 1024 \\
Lista 3 · Ejercicio 249 & AFN-λ & 12 & 4 & 2 & 4096 \\
Lista 3 · Ejercicio 250 & AFN-λ & 16 & 6 & 5 & 65536 \\
Lista 3 · Ejercicio 251 & AFN-λ & 18 & 4 & 4 & 262144 \\
Lista 3 · Ejercicio 252 & AFN-λ & 38 & 13 & 8 & 274877906944 \\
Lista 3 · Ejercicio 253 & AFN-λ & 44 & 22 & 17 & 17592186044416 \\
Lista 3 · Ejercicio 260 & AFN & 5 & 7 & 6 & 32 \\
Lista 3 · Ejercicio 261 & AFN & 5 & 5 & 3 & 32 \\
Lista 3 · Ejercicio 263 & AFN & 3 & 3 & 3 & 8 \\
\end{longtable}
